\documentclass[10pt,a4paper]{amsart}
\usepackage{amsmath,amssymb,amsthm,graphicx}
\usepackage{amsmath}
\usepackage{amssymb}
\usepackage{hyperref}
\newtheorem{Theorem}{Theorem}[section]
\newtheorem{Corollary}[Theorem]{Corollary}
\newtheorem{Lemma}[Theorem]{Lemma}
\newtheorem{Proposition}[Theorem]{Proposition}
{\theoremstyle{definition}
\newtheorem{Note}[Theorem]{Note}
\newtheorem{Definition}[Theorem]{Definition}
\newtheorem{Problem}[Theorem]{Problem}
}
\usepackage{cleveref}
\crefname{originaldefinition}{original definition}{original definitions}
\Crefname{originaldefinition}{Original definition}{Original definitions}
\makeatletter
\newlabel{def:ui}{{1.23}{source}{}{sourcepointer-1}{}}
\newlabel{def:ui@cref}{{[originaldefinition][0][]1.23}{[1][0][]source}{}{}{}}
\makeatother
\crefname{originallemma}{original lemma}{original lemmas}
\Crefname{originallemma}{Original lemma}{Original lemmas}
\makeatletter
\newlabel{lem:almostcom}{{1.25}{source}{}{sourcepointer-2}{}}
\newlabel{lem:almostcom@cref}{{[originallemma][0][]1.25}{[1][0][]source}{}{}{}}
\makeatother
\crefname{originaltheorem}{original theorem}{original theorems}
\Crefname{originaltheorem}{Original theorem}{Original theorems}
\makeatletter
\newlabel{thm:basis}{{1.21}{source}{}{sourcepointer-3}{}}
\newlabel{thm:basis@cref}{{[originaltheorem][0][]1.21}{[1][0][]source}{}{}{}}
\makeatother
\title{The monomials of degree at most n form a basis of the n-th filtration layer of the universal Askey-Wilson algebra}
\author{Paul Terwilliger}
\date{doi:10.3842/SIGMA.2011.069; CC BY4.0}
\begin{document}
\maketitle
\noindent\textit{Editorial note.} All delivered mathematical text is the author's own, copied from the pinned publisher source sources/sigma-2011-069-askey-wilson/pinned-original-source.tex; no statement or proof step was rewritten, reordered or removed.\par\medskip
\noindent\textit{Editorial note.} Preamble provenance: the retained preamble is the source's own layout adapter, if any, plus the author's macro block at lines 10-18 of sources/sigma-2011-069-askey-wilson/pinned-original-source.tex (everything between the publisher class metadata and \textbackslash{}begin\{document\}, comments excluded). Nothing was added, no macro was invented, and the publisher document class is deliberately not reproduced; the excerpt serves only as preamble, never as a span coordinate.\par\medskip
\noindent\textit{Editorial note.} External references are navigation pointers only; every deliberately kept step of the author's argument stays inside the counted spans. The displayed numbers of those pointers are the author's own numbering under the author's \textbackslash{}newtheorem declarations (the source has no subsections, so numbering is section.index), computed mechanically by the lane's scanner.\par\medskip
\noindent\textit{Editorial note.} Three references to other labelled items occur inside the counted spans (Definition def:ui, Lemma lem:almostcom, Theorem thm:basis); each is recorded once as a navigation stub with the matching kind and the author's own number, and every step that consumes them stays inside the counted span.\par\medskip
\noindent\textit{Editorial note.} The counted spans contain no bibliography citation, so no bibliography excerpt is attached to this card.\par\medskip
\noindent\textit{Allowed standard prerequisites (editorial).}\begin{itemize}
\item Linear independence of the monomials in the universal Askey-Wilson algebra. Same source, Theorem labelled thm:basis; invoked inside the counted proof at lines 1361-1363. Gives the independence half of the statement outright.
\item The filtration whose n-th layer is spanned by words of length at most n. Same source, Definition labelled def:ui; invoked inside the counted proof at lines 1365-1376. Supplies the spanning set by words in the generators, including the base case of the induction.
\item Swapping two adjacent letters modulo the previous layer. Same source, Lemma labelled lem:almostcom; invoked inside the counted proof at lines 1385-1392. Lets any word be reordered modulo Delta\_\{n-1\}, producing the monomial form needed for the spanning claim.
\end{itemize}\medskip
\ifdefined\hypertarget\hypertarget{sourcepointer-1}{}\fi
\noindent\textit{Source reference.} Original definition 1.23: sources/sigma-2011-069-askey-wilson/pinned-original-source.tex, source lines 1366--1366. Navigation pointer for \textbackslash{}ref\{def:ui\} used inside the counted span (line 1366); the argument that consumes it is retained in the span.\par\medskip
\ifdefined\hypertarget\hypertarget{sourcepointer-2}{}\fi
\noindent\textit{Source reference.} Original lemma 1.25: sources/sigma-2011-069-askey-wilson/pinned-original-source.tex, source lines 1385--1385. Navigation pointer for \textbackslash{}ref\{lem:almostcom\} used inside the counted span (line 1385); the argument that consumes it is retained in the span.\par\medskip
\ifdefined\hypertarget\hypertarget{sourcepointer-3}{}\fi
\noindent\textit{Source reference.} Original theorem 1.21: sources/sigma-2011-069-askey-wilson/pinned-original-source.tex, source lines 1363--1363. Navigation pointer for \textbackslash{}ref\{thm:basis\} used inside the counted span (line 1363); the argument that consumes it is retained in the span.\par\medskip
\begin{Theorem}\label{lem:unbasis}
For all integers $n\geq 0$ the following is
a basis for the~$\mathbb F$-vector space~$\Delta_n$:
\begin{gather}
\label{eq:basislist}
A^iB^jC^k\alpha^r \beta^s \gamma^t, \qquad
i,j,k,r,s,t\geq 0, \qquad
i+j+k+r+s+t \leq n.
\end{gather}
\end{Theorem}

\begin{proof}
The elements~(\ref{eq:basislist}) are linearly independent by
Theorem
\ref{thm:basis}. We show that the
elements~(\ref{eq:basislist}) span $\Delta_n$.
We will use induction on $n$. Assume $n\geq 2$; otherwise
the result holds by Def\/inition~\ref{def:ui}.
By Def\/inition~\ref{def:ui} $\Delta_n$ is
spanned by the set of elements of the form
$x_1x_2\cdots x_n$
where $x_i \in \lbrace 1, A,B,C,\alpha,\beta,\gamma\rbrace$
for $1 \leq i \leq n$.
Therefore $\Delta_n$ is spanned by the set of elements
of the form
$x_1x_2\cdots x_m$ where $0 \leq m\leq n$ and
$x_i \in \lbrace A,B,C,\alpha,\beta,\gamma\rbrace$
for $1 \leq i \leq m$.
Therefore $\Delta_n$ is spanned by $\Delta_{n-1}$ together with
the set of elements
of the form
$x_1x_2\cdots x_n$ where
$x_i \in \lbrace A,B,C,\alpha,\beta,\gamma\rbrace$
for $1 \leq i \leq n$.
Consider such an element
$x_1x_2\cdots x_n$.
By Lemma~\ref{lem:almostcom} and since
each of $\alpha$, $\beta$, $\gamma$ is central,
we f\/ind that
for $2 \leq j \leq n$,
\begin{gather*}
x_1\cdots  x_{j-1}x_j\cdots  x_n \in
\mathbb F x_1\cdots x_{j}x_{j-1}\cdots  x_n + \Delta_{n-1}.
\end{gather*}
By the above comments
$\Delta_n$ is spanned by $\Delta_{n-1}$ together with the
set
\begin{gather*}
A^iB^jC^k\alpha^r \beta^s \gamma^t, \qquad
i,j,k,r,s,t\geq 0, \qquad
i+j+k+r+s+t = n.
\end{gather*}
By this and induction
$\Delta_n$ is spanned by the elements~(\ref{eq:basislist}).
We have shown that
the elements~(\ref{eq:basislist}) form a basis  for the $\mathbb F$-vector space~$\Delta_n$.
\end{proof}

\begin{thebibliography}{99}

\footnotesize\itemsep=0pt
\bibitem{alpern}
Alperin R.C.,
Notes: ${\rm  {PSL}}_2(\mathbb Z)= \mathbb Z_2 \star \mathbb Z_3$,
\href{http://dx.doi.org/10.2307/2324963}{{\it Amer. Math. Monthly}} {\bf 100} (1993), 385--386.
\end{thebibliography}
\end{document}
